\documentclass[11pt]{article}
\usepackage[margin=1in]{geometry}
\usepackage[T1]{fontenc}
\usepackage{lmodern,microtype,hyperref,enumitem,amsmath,amssymb,xcolor,booktabs}
\usepackage[numbers,sort&compress]{natbib}
\hypersetup{colorlinks=true,urlcolor=blue,citecolor=blue,linkcolor=blue}
\title{Q19: What Are the Signs That AI Has Its Own Consciousness?}
\author{DLS 2026 --- module author: Kukla}
\date{October 2026}
\usepackage{url}
\begin{document}
\maketitle

\section*{Learning objectives}
\begin{itemize}[leftmargin=1.4em]
\item Distinguish consciousness from intelligence, capability, and self-report.
\item Summarize how leading scientific theories of consciousness have been applied to AI.
\item Reason carefully about why behavioral ``signs'' are weak evidence, and what better evidence looks like.
\end{itemize}

\begin{quote}\small\textbf{Companion videos (watch in order):}\begin{itemize}\setlength\itemsep{0pt}\item \emph{Why AI Isn't Going to Become Conscious} --- Anil Seth / TED (14:58). \url{https://www.youtube.com/watch?v=tJV-vdbZ388}\item \emph{How could we tell whether AI has become conscious?} --- Sabine Hossenfelder (10:23). \url{https://www.youtube.com/watch?v=CSTfgYynziw}\end{itemize}\end{quote}

\section{First, separate the concepts}
Student response \#01 asks what the signs are that an AI has its own consciousness. The first discipline
is to \emph{not} conflate consciousness with things it is routinely confused for: \textbf{intelligence}
(solving hard problems), \textbf{capability} (what a system can do), \textbf{agency} (pursuing goals),
and \textbf{self-report} (saying ``I am conscious''). A system can have any of these without
phenomenal experience, and---at least in principle---experience could exist without high intelligence.
``Consciousness'' here means \emph{phenomenal} consciousness: that there is something it is like to be
the system \citep{nagel1974bat}. This is distinct from \emph{access} consciousness (information being
available for reasoning and report), which is more tractable but not the same question.

\section{Why self-report is nearly worthless as evidence}
An LLM trained on human text will fluently produce first-person claims of feelings because that is what
the training distribution contains---not because it has introspective access to inner states. The
report is \emph{caused by} the data and the prompt, not by an experience. Worse, these claims are
steerable: the same model will assert and deny consciousness depending on the prompt. So ``it said it
was conscious'' is close to zero evidence. This is the single most important point for a student to
internalize.

\section{Theory-driven indicators}
A more principled approach asks what each scientific theory of consciousness \emph{requires}, then checks
whether an architecture has it \citep{butlin2023consciousness}. Examples:
\begin{itemize}[leftmargin=1.4em]
\item \textbf{Global Workspace Theory}: a bottleneck that broadcasts selected information globally. One
can ask whether a model has a functional global workspace (some do, partially, by design).
\item \textbf{Higher-Order theories}: representations \emph{about} the system's own representations.
\item \textbf{Integrated Information Theory} (IIT): consciousness tracks a system's irreducible
integrated information $\Phi$ \citep{tononi2016iit}. Computing $\Phi$ for real networks is intractable,
and IIT makes the sharp, counterintuitive claim that a feed-forward network has $\Phi=0$ regardless of
behavior.
\item \textbf{Recurrent / predictive-processing} accounts: emphasize recurrence and self-modeling.
\end{itemize}
Butlin et al.\ \citeyearpar{butlin2023consciousness} turn these into a checklist of ``indicator
properties'' and conclude current systems satisfy few of them, but that there is no obvious technical
barrier to building systems that satisfy more.

\section{Why this is genuinely hard}
Three deep problems block a clean test. (1) The \textbf{other-minds problem}: we infer consciousness in
other humans by analogy from shared biology and behavior---an AI breaks that analogy. (2) \textbf{Gaming}:
a capable system can produce any behavioral ``sign'' we name, so behavioral tests are not robust to a
system optimized to pass them. (3) \textbf{Theory disagreement}: the theories above disagree with each
other, so an ``indicator'' is only as trustworthy as a theory we have not confirmed. The responsible
scientific posture is therefore probabilistic and theory-relative, not a yes/no verdict.

\section{A concrete walk through one theory}
To make ``theory-driven indicator'' concrete, take Global Workspace Theory (GWT). GWT holds that a
mental state is conscious when its content is \emph{globally broadcast}: selected by an attentional
bottleneck and made available to many downstream processes (memory, reasoning, report)
\citep{dehaene2014consciousness}. Applied to AI, the indicator question becomes architectural and
checkable: does the system have a limited-capacity workspace that gates information, and is the gated
content broadcast to otherwise independent modules? A vanilla transformer has attention (a selection
mechanism) but not a persistent capacity-limited workspace broadcasting to specialized subsystems, so it
scores \emph{partially} on GWT indicators---and one could deliberately build a system that scores higher.
Crucially, scoring high on GWT indicators would make the system ``conscious according to GWT,'' which is
only meaningful if GWT is correct---a question neuroscience has not settled. This is why serious work
reports indicator scores \emph{per theory} with explicit uncertainty, rather than a single verdict, and
why Rick's machine-consciousness program runs sham-controlled contrasts tied to specific theories (global
workspace, recurrence, higher-order, integrated-information proxies) instead of asking the model how it
feels.

\section{The moral stakes of getting it wrong}
Two asymmetric errors loom. A \emph{false positive}---attributing experience to a system that has none---
could paralyze useful work and misdirect moral concern. A \emph{false negative}---denying experience to a
system that has it---risks real moral harm at scale. Because the theories disagree and the evidence is
thin, the responsible stance treats the question as a live uncertainty with non-trivial probability mass
on several answers, and separates it cleanly from capability and welfare questions that can be addressed
independently. Students should leave able to hold ``we do not know'' without collapsing it into either
``obviously conscious'' or ``obviously just a program.''

\section{What better evidence would look like}
Convergent evidence across independent theories (a system scoring high on global-workspace,
higher-order, \emph{and} recurrence indicators), causal-mechanistic tests (interventions that
manipulate the putative substrate and change the putative markers), and architectures where the markers
were \emph{not} optimized for---so passing them is informative. Even then, the conclusion is
``consistent with'' not ``proves.'' This is an active research area; Rick's own machine-consciousness
program runs exactly these theory-grounded, sham-controlled contrasts rather than trusting self-report.

\section{Summary}
Consciousness is not intelligence, capability, agency, or self-report; self-report in LLMs is almost
worthless because it is a trained behavior. The scientific route is theory-driven indicator properties
(global workspace, higher-order, IIT, recurrence), but theories disagree, behavioral signs can be gamed,
and the other-minds problem is real. The honest answer: there is no single reliable ``sign,'' only
theory-relative, probabilistic, convergent evidence---and current systems satisfy few indicators.

\bibliographystyle{plainnat}
\bibliography{bibliography}
\end{document}
